% eksperymentalne tłumaczenie na włoski

%

\section{Stożkowe} % lang-pl

\section{Coniche} % lang-it

% TODO: Alhazen => https://en.wikipedia.org/wiki/Straightedge_and_compass_construction => niemożliwy
% https://en.wikipedia.org/wiki/Alhazen%27s_problem

\begin{problem}[bilardowy Alhazena] % lang-pl
    Dany jest okrąg $\Gamma$ oraz dwa punkty $A$, $B$, obydwa w jego wnętrzu lub obydwa na zewnątrz. % lang-pl
\begin{problem}[problema biliardistico di Alhazen] % lang-it
    
    Skonstruować trójkąt równoramienny wpisany w okrąg $\Gamma$ tak, żeby punkt $A$ leżał na jednym ramieniu, zaś $B$ na drugim. % lang-pl
    Sono dati un cerchio $\Gamma$ e due punti $A$, $B$, entrambi al suo interno oppure entrambi al suo esterno. % lang-it
    
    Costruire un triangolo isoscele inscritto nel cerchio $\Gamma$ in modo che il punto $A$ appartenga a un lato obliquo e il punto $B$ all'altro. % lang-it
\end{problem}

Tematyką zajmie się już Ptolemeusz około II wieku naszej ery, ale dopiero XI-wieczny matematyk arabski, Alhazen (Hasan Ibn al-Haytham), sformułuje problem ogólniej, a oprócz tego przedstawi też jego rozwiązanie w książce \emph{Kitab al-Manazur} (łacińskie \emph{De aspectibus}, o optyce). % lang-pl

Około 1965 roku, amerykański matematyk Jack Elkin rozwiąże problem algebraicznie; pojawi się tam pierwiastek równania czwartego stopnia, więc trójkąta nie da się wykreślić samym cyrklem z linijką. % lang-pl
L'argomento sarà studiato già da Tolomeo intorno al II secolo d.C., ma soltanto il matematico arabo dell'XI secolo Alhazen (Hasan Ibn al-Haytham) ne formulerà una versione più generale e ne presenterà una soluzione nel \emph{Kitab al-Manazur} (in latino \emph{De aspectibus}), dedicato all'ottica. % lang-it

Al-Haytham i inni (na przykład Christiaan Huygens) wykorzystają wcześniej przecięcie stożkowych. % lang-pl
Intorno al 1965 il matematico statunitense Jack Elkin risolverà il problema algebricamente; comparirà una radice di un'equazione di quarto grado, per cui il triangolo non può essere costruito con la sola riga e il compasso. % lang-it
Al-Haytham e altri (per esempio Christiaan Huygens) utilizzeranno invece l'intersezione di coniche. % lang-it


Problem ma prostą interpretację fizyczną, dlatego też znalazł się w dziele o optyce. % lang-pl
Il problema ha una semplice interpretazione fisica, motivo per cui compare in un'opera di ottica. % lang-it

Promień światłą biegnie z punktu $A$ aż uderzy w powierzchnię kuli, po czym zmienia kierunek tak, że kąty między promieniem przed i po uderzeniu oraz normalną do powierzchni sfery są równe. % lang-pl
Un raggio luminoso parte dal punto $A$, colpisce la superficie di una sfera e cambia direzione in modo che gli angoli formati dal raggio incidente e riflesso con la normale alla superficie siano uguali. % lang-it

Szukamy takiego punktu, by po odbiciu trafić w punkt $B$. % lang-pl
Cerchiamo il punto di riflessione tale che il raggio raggiunga il punto $B$ dopo il rimbalzo. % lang-it


Więcej przeczytać można u Hartshorne'a \cite[s. 278]{hartshorne2000}. % lang-pl
Per ulteriori dettagli si veda Hartshorne \cite[p. 278]{hartshorne2000}. % lang-it

https://www.deltami.edu.pl/2017/06/zadanie-alhazena/

%